\section{Related Work}
\label{sec:related}

\paragraph{Accident anticipation and the benchmark tradition.}
Accident anticipation has moved from early temporal modeling for DAD~\citep{ref4} to larger multi-branch and graph-based predictors~\citep{ref45,ref17,ref32,ref36,ref37,ref208}. Large-capacity fusion systems now report strong AP and mTTA values on benchmark splits~\citep{liao2024crash,liu2025ccaf,ZHANG2025103173}, and several lines also add auxiliary interpretability signals~\citep{ref2,ref102}. The field therefore has a mature predictive toolbox, but most methods still optimize toward one central axis: score quality over time. In safety-critical forecasting, this leaves an unclosed question about what the model is using, and whether that ``why'' can be governed across runs and deployment settings.
The broader interpretability literature frames this same gap as one between local
faithfulness and reusable trust, with recent work emphasizing explicit interfaces
 explanations in risky settings~\citep{ribeiro2016loco,lundberg2017unified,doshi2017towards,rudin2019stop}. Some policy-oriented lines for safety timing introduce explicit uncertainty and control objectives~\citep{osband2016bootstrapped,konda1999actor}, yet they optimize decision behavior rather than providing a governed semantic vocabulary as an explicit method object.
In road-safety settings, deployment practice also raises this question to system
validity: interpretation is part of validation, not only visualization~\citep{who2023global,NHTSA2023}.
Related forecasting and trajectory work motivates the same concern about whether the
underlying representation is trustworthy under operational decisions~\citep{liao2024cognitive,ref13,ref23}.
The transportation literature further shows that operator safety judgment is tied to
semantic alignment and interpretable cueing, not solely raw confidence scores~\citep{kim2025understanding,guan2025experimental}.
Recent works on richer localization and accident context estimation~\citep{ref100,wang2024deepaccident} strengthen practical utility, while still leaving the semantic interface design itself mostly implicit.
Related datasets and precursor lines, including uncertainty-aware anticipation
and driver-attention-oriented benchmarks~\citep{ref1,ref33,fang2021dada,karim2023attention},
similarly emphasize predictive behavior more than governed semantic interfaces.

\paragraph{Governance and human factors context.}
Recent safety standards and practitioner guidance stress that model evidence is one part of readiness.
Validation protocols for safety-critical forecasting increasingly require interpretable
artifacts, explicit uncertainty behavior, and auditability paths, especially when
warnings trigger human or automation actions~\citep{who2023global,NHTSA2023}.
\method{} positions the semantic interface as one such artifact: a method-level
object whose names, families, and timing behavior are reviewable across protocol
choices.

In high-stakes settings, this matters because interpretability without governance
is often non-actionable. Concept-level claims require evidence about how concepts are
constructed, maintained, and constrained across runs; the paper therefore treats the
ontology and scoring pipeline as part of the model method and not as a detached
annotation step. This aligns with the broader call in high-stakes ML to keep
modeling choices auditable by design rather than evaluated only after training
~\citep{doshi2017towards,rudin2019stop}.

\paragraph{Post-hoc explanation versus embedded semantics.}
Post-hoc methods such as feature attribution and verbalization provide useful
diagnostics (e.g., LIME and SHAP) but are often attached after prediction.
DRIVE adds attention
visualizations to the output path~\citep{ref2}; W3AL introduces natural language
reporting around location and timing~\citep{liao2024and}. These designs increase
transparency to users but do not usually enforce a stable semantic state that is
shared by training, evaluation, and intervention. In our setting, this distinction is
critical: an attached explanation can show what looked important for one
prediction, while an embedded semantic interface must support reusable claims
about what risk concepts are consistently represented by the model.
\method{} positions itself in this gap by making the concept state part of the
model's forward computation.

\paragraph{Interpretability and concept-level learning.}
Concept bottlenecks provide a principled intermediate layer that remains
human-interpretable and mathematically tractable~\citep{tishby2000information,koh2020concept}.  
Subsequent work has reduced reliance on dense annotations and expanded
concept-supervision design~\citep{zarlenga2022concept,yuksekgonul2023posthoc,oikarinen2023label}.  
In accident anticipation, those ideas are harder to apply because concepts must
be temporally meaningful and task-governable across long horizons. This is why
our method treats concept vocabulary construction as a first-class component,
not a fixed side input.

This framing also distinguishes \method{} from post-hoc explanation pipelines.
In the broader accident-anticipation literature, most prior work does not treat
concept vocabulary as a tunable, auditable method component, so semantic signals
cannot be directly governed across matched launches.
Attribution methods such as Grad-CAM
\citep{selvaraju2017gradcam} explain individual predictions, while
concept-activation techniques such as TCAV
\citep{kim2018tcav} test global class-level concept sensitivity.
Both are useful, but neither becomes a reusable interface artifact unless the
concept layer is constructed, constrained, and audited as part of the model.
This also makes the actor-policy caution more central: prior reinforcement-learning
and credit-assignment work emphasize that policy- or intervention-level claims require explicit control-focused validation, not prediction-only benchmarking~\citep{konda1999actor,meulemans2023cocoa}.
Our contribution is therefore not only semantic readability, but the governance
of that interface under matched launchers and matched family-level controls.
Automatic concept discovery lines such as ACE~\citep{ghorbani2019ace} are also
relevant here: they motivate scalable concept proposal, but our focus is the
downstream governance and deployment of a fixed, auditable concept interface.

\paragraph{Language-grounded ontology for multimodal safety systems.}
Recent progress in vision-language pretraining enables language to act as
semantic glue for risky-scene understanding~\citep{radford2021learning,jia2021align}.  
In driving, VL supervision has improved scene grounding and text-aware
reasoning~\citep{liao2024gpt,shi2025scvlm}. In contrast to trajectory-prediction systems that rely heavily on map constraints and route semantics~\citep{ref215}, we use it as a conservative
bootstrap: CLIP-style pseudo-labeling and paraphrase stress tests produce
candidate semantic structure while avoiding claims of full semantic correctness.
Temporal distillation then anchors the interface to future-aware representations,
so the semantic layer is coupled to anticipation timing~\citep{fang2022traffic}.

\paragraph{Ontology governance as method, not ornament.}
Compared with systems where semantic terms are implicit prompts, prompt templates,
or uncontrolled phrase pools, our pipeline explicitly performs mining,
canonicalization, family balancing, and light review. This changes the research
question from \emph{``Can we predict with higher AP?''} to \emph{``How does
semantic design change behavior under matched training and evaluation?''}.  
That shift aligns with current language-oriented evaluation culture because the
interface itself is now an accountable artifact: its names, merges, and coverage can
be inspected and compared.

\paragraph{Position in this paper.}
Our paper is not attempting to turn accident anticipation into a text-generation
task. Instead, it studies whether a governed, language-grounded interface can be
designed for a forecasting pipeline, and whether that interface changes what can
be audited and compared. In this framing, \method{} contributes primarily at the
representation-design layer: a reproducible ontology construction path and a
matched-evaluation protocol that turns ontology choice into an explicit variable.
\method{} is strongest where this differs from prior accident systems: it links
concept construction, model training, and intervention-oriented auditing into one
methodological unit rather than treating semantics as a visualization overlay.
